\section{Complete canonical correlation intervals and a sharp success-probability frontier}
\label{sec:probability-frontier}

The centre and radial-width bound determines an entire attainable interval,
not only its upper endpoint. A separate operational refinement keeps track of
the probability of the filtering event. Throughout this section a canonical
input has Bloch data $(c\nn,0,S)$, where $\norm\nn=1$ and $0\leq c\leq1$.
When $c=0$, the direction $\nn$ is a specified reference direction; it is not
claimed to be determined by the vanishing centre. Set
\[
 t=\norm{S^{\mathsf T}\nn},\qquad \tau=\norm S_F^2.
\]
All Kraus probabilities below refer to the canonical \emph{input} state.

\begin{theorem}[Complete canonical correlation interval]\label{thm:canonical-interval-v2}
The pair $(c,t)$ is feasible exactly when $0\leq t\leq1-c$. At every feasible
pair, the complete set of squared correlation norms is
\begin{equation}\label{eq:canonical-full-interval-v2}
 \boxed{\quad t^2\leq\tau\leq1-c^2+2t^2,\quad}
\end{equation}
with every value in this closed interval attained by a density matrix.
Both endpoints and all singular parameter cases are included.

The lower endpoint is attained exactly when $S=\nn\mathbf u^{\mathsf T}$ for
some $\norm{\mathbf u}=t$. Thus the lower-endpoint state is unique up to local
unitaries. The upper-endpoint state is also unique up to local unitaries and
has the matrix form displayed below. Its spectrum is
$\{(1+t)/2,(1-t)/2,0,0\}$.
\end{theorem}

\begin{proof}
A local rotation puts $\nn=\mathbf e_z$; a Bob rotation puts
$S^{\mathsf T}\nn=t\mathbf e_z$. Positivity of the conditional qubit with
Bob direction $\mathbf e_z$ gives $c+t\leq1$. The lower bound follows from
\[
 \tau-t^2=\norm{(\Id-\nn\nn^{\mathsf T})S}_F^2\geq0,
\]
and equality is precisely the stated rank-one condition. Its diagonal
representative is physical for exactly $t\leq1-c$.
For the upper bound, let
\[
 U=(\nn\cdot\bs)\sigma_y\otimes\sigma_y,
 \qquad \sigma_{\nn}=U\overline\sigma U^\dagger.
\]
Both $\sigma$ and $\sigma_{\nn}$ are positive semidefinite, so their trace
product is nonnegative. Pauli orthogonality gives
\[
 4\Tr(\sigma\sigma_{\nn})=1-c^2+2t^2-\tau\geq0.
\]
This argument does not require either state to be invertible.

For explicit attainability define
\[
 a=\frac{1+c+t}{4},\quad b=\frac{1+c-t}{4},\quad
 e=\frac{1-c-t}{4},\quad d=\frac{1-c+t}{4}.
\]
All four numbers are nonnegative. For $0\leq\alpha\leq1$ set
\begin{equation}\label{eq:interval-family-v2}
 \sigma_{c,t,\alpha}=\begin{pmatrix}
 a&0&0&\alpha\sqrt{ad}\\
 0&b&\alpha\sqrt{be}&0\\
 0&\alpha\sqrt{be}&e&0\\
 \alpha\sqrt{ad}&0&0&d
 \end{pmatrix}.
\end{equation}
Its two $2\times2$ blocks have determinants $(1-\alpha^2)ad$ and
$(1-\alpha^2)be$. Hence it is a density matrix, including when a diagonal
entry vanishes. Its Bloch data are
\[
 \cc=c\mathbf e_z,\quad\bb=0,\quad
 S=\operatorname{diag}\bigl(2\alpha(\sqrt{ad}+\sqrt{be}),
                    2\alpha(\sqrt{be}-\sqrt{ad}),t\bigr),
\]
so
\begin{equation}\label{eq:interval-param-v2}
 \tau=t^2+\alpha^2(1-c^2+t^2).
\end{equation}
For $1-c^2+t^2>0$, any requested $\tau$ in
\eqref{eq:canonical-full-interval-v2} is obtained by choosing
$\alpha=\sqrt{(\tau-t^2)/(1-c^2+t^2)}$. If that denominator vanishes,
$c=1,t=0$ and the interval consists of the single value zero.

For completeness, consider upper-endpoint rigidity. In the chosen axes,
every canonical state has block form
\[
 \sigma=\begin{pmatrix}A&C\\C^\dagger&D\end{pmatrix},\qquad
 A=\operatorname{diag}(a,b),\quad D=\operatorname{diag}(e,d),\quad\Tr C=0.
\]
If $0<t<1-c$, both diagonal blocks are positive definite. Block positivity
is equivalent to $C=A^{1/2}KD^{1/2}$ with $K$ a contraction. The matrix
$(|K_{ij}|^2)$ has row and column sums at most one. The elementary
rearrangement inequality for such a $2\times2$ matrix gives
\[
 \norm C_F^2\leq ad+be.
\]
One way to verify the equality condition is to extend the matrix of squared
moduli to a doubly stochastic matrix. Because all four weights are strictly
positive, equality permits no added entry. Every doubly stochastic
$2\times2$ matrix is a convex combination of the two permutation matrices;
the difference between their weighted sums is
$(a-b)(d-e)=t^2/4>0$. Equality therefore forces $K$ to be antidiagonal with
both nonzero entries of modulus one. Alice and Bob rotations about $z$
remove its two phases, giving \eqref{eq:interval-family-v2} with $\alpha=1$.

If $t=0$ and $c<1$, the two blocks are scalar matrices. Equality gives
$C=\sqrt{ad}\,V$, where $V$ is a traceless unitary. Every traceless
$2\times2$ unitary is a scalar phase times a traceless Hermitian unitary;
Bob conjugation and an Alice phase rotation reduce it to $\sigma_x$.
This again gives the displayed upper-endpoint representative. If
$t=1-c>0$, the diagonal entry $e$ vanishes. Positivity forces its entire row
and column to vanish; the remaining coherence has modulus at most
$\sqrt{ad}$, and upper equality fixes that modulus. The cases $(c,t)=(0,1)$
and $(1,0)$ follow directly from the same zero-row rule. Thus no inverse of
a zero diagonal block was used. At $\alpha=1$, the two blocks have traces
$(1+t)/2$ and $(1-t)/2$ and determinant zero, proving the eigenvalue claim.
\end{proof}

\paragraph{Ranks in the explicit interpolation.}
For $c<1$ and $t<1-c$, the state \eqref{eq:interval-family-v2} has rank four
for $\alpha<1$ and rank two for $\alpha=1$. For $0<c<1$ and $t=1-c$, these
ranks are three and two. At $(c,t)=(0,1)$ they are two and one. At
$(c,t)=(1,0)$ the matrix is the fixed rank-two product state
$|0\rangle\langle0|\otimes\Id/2$, for every $\alpha$.

\begin{lemma}[Exact filtering probability]\label{lem:success-probability-v2}
Let $\sigma_B=\Id/2$ and let $G$ be a nonzero-probability Bob Kraus operator,
$G^\dagger G\preceq\Id$. Its success probability is
\[
 p=\Tr[(\Id\otimes G)\sigma(\Id\otimes G^\dagger)]
   =\tfrac12\Tr(G^\dagger G)>0.
\]
Writing $F=G/\sqrt p$ gives the unique normalized effect representation
\[
 F^\dagger F=\Id+\mathbf x\cdot\bs,\qquad r=\norm{\mathbf x}\leq1.
\]
For a fixed normalized effect, the exact maximal success probability is
\begin{equation}\label{eq:pmax-v2}
 p_{\max}(r)=\frac1{1+r}.
\end{equation}
Every probability $0<p\leq p_{\max}(r)$ is attained by
$G=\sqrt p\,F$. Its normalized output is exactly
$(\Id\otimes F)\sigma(\Id\otimes F^\dagger)$.
\end{lemma}

\begin{proof}
The trace identity follows by tracing out Alice and using $\sigma_B=\Id/2$.
The positive matrix $F^\dagger F$ has trace two and therefore has the indicated
Bloch representation, with eigenvalues $1\pm r$. The Kraus condition is
\[
 \Id-pF^\dagger F\succeq0
 \quad\Longleftrightarrow\quad p(1+r)\leq1.
\]
This proves both necessity and attainability, including $r=1$. The scalar
factor $p$ cancels from the normalized output. If $p>1/2$, the inequality
forces $r<1$, so $G$ is invertible. Conversely a singular nonzero qubit Kraus
operator has $r=1$ and success at most $1/2$.
\end{proof}

For a required success probability $p_0\in[1/2,1]$, define
\begin{align}
 R(p_0)&=\frac{1-p_0}{p_0},&
 D(c)&=2-2c-c^2,\\
 H_c(r)&=(1-r^2)(3-4c)+[c+(1-c)r]^2,\label{eq:prob-H-v2}\\
 r_*(c)&=\frac{c(1-c)}{D(c)}\quad(0\leq c\leq2/3),&
 \widehat r(c,p_0)&=
 \begin{cases}\min\{R(p_0),r_*(c)\},&c\leq2/3,\\
 R(p_0),&c>2/3.
 \end{cases}\label{eq:prob-ropt-v2}
\end{align}
The definition of $r_*$ at $c=2/3$ is regular and gives $r_*=1$.

\begin{theorem}[Sharp probability-constrained correlation envelope]
\label{thm:probability-envelope-v2}
Over all canonical input density matrices with centre length $c\in[0,1]$
and all Bob Kraus operators of success probability at least
$p_0\in[1/2,1]$, the exact greatest squared correlation norm is
\begin{equation}\label{eq:prob-envelope-v2}
 \boxed{\quad \mathcal G(c,p_0)=H_c(\widehat r(c,p_0)).\quad}
\end{equation}
This is a maximum, including all endpoints. An explicit maximizer is
\begin{align}
 \sigma_c&=\frac c2|00\rangle\langle00|
 +\frac12(|01\rangle+\sqrt{1-c}|10\rangle)
           (\langle01|+\sqrt{1-c}\langle10|),\label{eq:prob-obese-v2}\\
 G_{c,p_0}&=\sqrt{p_0}\,
     \operatorname{diag}\bigl(\sqrt{1-\widehat r},
                              \sqrt{1+\widehat r}\bigr).
\end{align}
The actual success probability of this witness is exactly $p_0$.
\end{theorem}

\begin{proof}
Normalize any admissible Kraus operation as in
Lemma~\ref{lem:success-probability-v2}. If its actual success is $p\geq p_0$,
then $p_0(1+r)\leq p(1+r)\leq1$, hence $0\leq r\leq R(p_0)$.
The exact correlation identity and positivity give
\[
 F_3^2=(1-r^2)(\tau-c^2)+\norm{\cc+S\mathbf x}^2
 \leq(1-r^2)(3-4c)+[c+(1-c)r]^2=H_c(r).
\]
The canonical bound used here is $\tau\leq3-4c+c^2$. The marginal bound
follows by writing
$\cc+S\mathbf x=(1-r)\cc+r(\cc+S\mathbf u)$, with
$\mathbf u=\mathbf x/r$ when $r>0$; the conditional endpoint
$\cc+S\mathbf u$ lies in the Bloch ball. The case $r=0$ follows directly.
The identity and inequalities hold also for $r=1$ by direct substitution.

Expanding gives
\[
 H_c(r)=3-4c+c^2+2c(1-c)r-D(c)r^2.
\]
For $c\leq2/3$, $D(c)>0$ and the quadratic reaches its maximum at $r_*(c)$;
clipping to $[0,R]$ gives \eqref{eq:prob-ropt-v2}. For $c\geq2/3$, the
quadratic is nondecreasing on $[0,1]$. Indeed if $D\geq0$, its derivative
is bounded below by $2[c(1-c)-D]=2(3c-2)\geq0$, and if $D<0$ its derivative
is nonnegative directly. The endpoint $R$ is therefore optimal.

The state \eqref{eq:prob-obese-v2} is a positive sum of two rank-one
operators, has trace one, has Bob marginal $\Id/2$, and has centre length
$c$. Direct multiplication by the displayed diagonal filter gives
$F_3^2=H_c(\widehat r)$. Its normalized effect has radius $\widehat r$;
$p_0(1+\widehat r)\leq1$ proves the actual Kraus condition. Its success is
$p_0$ by the canonical trace identity. This proves the maximum over the
full physical class, not only over a scalar relaxation.
\end{proof}

\begin{corollary}[Exact activation frontier]\label{cor:success-threshold-v2}
There exists a canonical input of centre length $c$ and a Bob Kraus operation
of success probability at least $p_0\in[1/2,1]$ whose output has $F_3>1$
if and only if
\begin{equation}\label{eq:success-critical-v2}
 \boxed{\qquad c<C_{\rm crit}(p_0):=
 \frac{2}{1+p_0+\sqrt{(1-p_0)^2+2}}.\qquad}
\end{equation}
The equality curve is nonviolating. It interpolates between
$C_{\rm crit}(1/2)=2/3$ and $C_{\rm crit}(1)=2-\sqrt2$.
\end{corollary}

\begin{proof}
Let $E(c)=2-4c+c^2$. The exact factorization is
\begin{equation}\label{eq:prob-factor-v2}
 H_c(r)-1=(1-r)[D(c)r+E(c)].
\end{equation}
If there is a violating point with $r\leq R$, then $c<2/3$ by the global
nonviolation bound, so $D(c)>0$. The linear second factor must be positive
and is increasing in $r$. It follows that
\[
 g(c,p_0):=p_0[D(c)R(p_0)+E(c)]
          =2-2(1+p_0)c+(2p_0-1)c^2>0.
\]
For $p_0>1/2$, the converse follows by evaluating the endpoint:
\[
 H_c(R)-1=\frac{(2p_0-1)g(c,p_0)}{p_0^2}.
\]
At $p_0=1/2$, this endpoint identity is identically zero and cannot establish
activation. Instead $g=2-3c>0$ means $c<2/3$; the interior vertex
$r_*(c)<1=R$ gives $H_c(r_*)=G(c)>1$.

Put $s=\sqrt{(1-p_0)^2+2}$. Then
\[
 2g(c,p_0)=[2-(1+p_0+s)c][2-(1+p_0-s)c].
\]
For $0\leq c\leq1$ and $1/2\leq p_0\leq1$, the second factor is strictly
positive: $s\geq1$, hence $(1+p_0-s)c\leq1$. The sign of $g$ is therefore
exactly the sign of the first factor, proving the claimed strict threshold.
\end{proof}

\paragraph{Nonbinding budgets and equality for $p_0>1/2$.}
For $0\leq c<2/3$, the probability constraint leaves the unrestricted maximum
unchanged exactly when
\begin{equation}\label{eq:nonbinding-prob-v2}
 p_0\leq p_*(c):=\frac1{1+r_*(c)}
                  =\frac{D(c)}{2-c-2c^2}.
\end{equation}
Above that value the maximum is the endpoint expression $H_c(R(p_0))$.
For $2-\sqrt2<c<2/3$, the largest allowed success floor for strict activation
is equivalently
\[
 p_0<\frac{D(c)}{2c(1-c)}.
\]
At $c=2-\sqrt2$, activation occurs for every $p_0<1$ and fails at $p_0=1$.
For $c<2-\sqrt2$, the displayed canonical input
\eqref{eq:prob-obese-v2} itself already violates.

For $0<c<1$ and $p_0>1/2$, an extremizer must saturate the sharp joint
centre/radius bound at its uniquely optimal radius $\widehat r<1$.
The canonical input is therefore a local-unitary copy of
\eqref{eq:prob-obese-v2}. When $\widehat r>0$, its normalized effect points
toward the pure contact of that canonical ellipsoid; in the displayed axes
it is $\Id-\widehat r\sigma_z$. A post-filter Bob unitary is unrestricted.
When $\widehat r=0$, the normalized effect is the identity. In a strictly
binding case $\widehat r=R$, the actual success of every maximizer is exactly
$p_0$. In a nonbinding case $\widehat r=r_*$, the actual success can be any
$q\in[p_0,p_*(c)]$, by scalar rescaling of the same normalized filter.
At the crossover these two descriptions agree. For $c=0$, the canonical
maximizers are maximally entangled pure states and the optimal normalized
effect is the identity. For $c=1$, the canonical state is
$a\otimes\Id/2$ with $a$ pure, and $F_3^2=r^2$; any filter direction is
allowed, with $r=R$.

\paragraph{Singular endpoints and the meaning of the centre constraint.}
At $p_0=1/2$, all radii through $r=1$ are feasible. For $c<2/3$, the optimum
still uses the invertible radius $r_*<1$. For $2/3\leq c\leq1$, the maximum
is one and is attained at $r=1$ by a singular Kraus operator selecting a pure
contact point. More generally, among canonical inputs in this range,
a singular filter attains one precisely when its selected conditional
Alice state is pure; the normalized output is then a pure product.
Many canonical input orbits with pure ellipsoid contacts can therefore
attain this singular maximum. No unique obese-input classification is
claimed there.

For $p_0>1/2$, all admissible filters are invertible and preserve the steering
ellipsoid and its centre. At $p_0=1/2$, the singular maximizing output has
centre length one, even if the specified \emph{input} centre is $c<1$.
This does not contradict the nonattainment of the global fixed-\emph{output}
centre envelope for $2/3\leq c<1$. If the success floor is instead
$0<p_0<1/2$, every normalized radius in $[0,1]$ is feasible; the maximum is
$G(c)$, interpreted as an attained value one on the large-centre branch,
and the activation threshold remains $c<2/3$.

\paragraph{Exact formalization boundary for this section.}
The complete physical interval, including its interpolation witness and
lower-endpoint equality condition, is checked in
\path{MilneGeometryInterval.lean}, principally by
\path{geometryPhysicalSkewSpectrum_eq_interval} and
\path{geometry_radial_lower_equality_iff}.
The upper-endpoint local-unitary classification and the listed interpolation
ranks are written proofs in this section.

\path{MilneKrausProbability.lean} checks actual complex Kraus positivity,
normalization, success probability, and output density. Its principal
lemmas are \path{scaledKraus_admissible_iff},
\path{fixed_filter_maximal_success},
\path{canonicalKrausProbability_is_actual},
\path{normalizedKrausOutput_density}, and
\path{canonical_complex_joint_bound}.
\path{MilneProbabilityScalar.lean} checks the clipped optimizer and exact
activation curve, including the separate $p_0=1/2$ argument, through
\path{successOptimalRadius_maximizes} and
\path{successEnvelope_violation_iff_center}.
\path{MilneProbabilityEnvelope.lean} states the optimization over
\emph{actual} density matrices and contracting complex Kraus operators.
Its main theorems are \path{successPhysicalSpectrum_isGreatest} and
\path{successPhysicalSpectrum_violation_iff}.
The exact saturation below probability one half is also checked by
\path{successPhysicalSpectrum_low_probability} and
\path{looseSuccessPhysicalSpectrum_isGreatest}; singularity exclusion above
one half is checked by
\path{KrausAdmissible_det_ne_zero_of_success_gt_half}.
The additional local-unitary equality classification and probability
nonbinding descriptions are written consequences; they are not represented
as separately checked Lean theorems here.
